تابع
$f(x,y)=50x^2+y^2$
یک تابع درجه‌دو محدب با انحنای متفاوت در دو راستا است. ژاکوبین/گرادیان و هسین آن:
\begin{flalign*}
\nabla f(x,y) &= 
\begin{bmatrix}
\frac{\partial f}{\partial x}\\[2pt]
\frac{\partial f}{\partial y}
\end{bmatrix}
=
\begin{bmatrix}
100x\\
2y
\end{bmatrix},
\qquad
\nabla^2 f(x,y)=
\begin{bmatrix}
100 & 0\\
0 & 2
\end{bmatrix}.
&&
\end{flalign*}

\begin{itemize}
\item \lr{Gradient Descent}
\begin{itemize}
\item[(الف)]
\begin{flalign*}
\nabla f(x,y)=\begin{bmatrix}100x\\2y\end{bmatrix}. &&
\end{flalign*}

\item[(ب)]
چون هسین قطری است، بیشترین انحنا در راستای $x$ برابر $100$ و در راستای $y$ برابر $2$ است (مسئله بدشرط با نسبت $100/2=50$).
برای پایداری \lr{GD} روی تابع درجه‌دو، نرخ یادگیری ثابت باید با بزرگ‌ترین مقدار ویژه سازگار باشد:
\begin{flalign*}
0<\eta<\frac{2}{L},\qquad L=\lambda_{\max}(\nabla^2 f)=100 \;\;\Rightarrow\;\; \eta<0.02. &&
\end{flalign*}
اگر $\eta$ را به‌اندازه‌ای بزرگ بگیریم که در راستای $y$ سریع باشد، در راستای $x$ نوسان/واگرایی رخ می‌دهد؛ اگر هم $\eta$ را کوچک کنیم تا $x$ پایدار شود، حرکت در راستای $y$ بسیار کند می‌شود.

\item[(ج)]
با $\eta=0.1$ و $(x_0,y_0)=(1,1)$:
\begin{flalign*}
\nabla f(x_0,y_0) &= (100,2), \\
x_1 &= x_0-\eta(100x_0)=1-0.1(100)= -9, \\
y_1 &= y_0-\eta(2y_0)=1-0.1(2)=0.8. &&
\end{flalign*}
تکرار دوم:
\begin{flalign*}
\nabla f(x_1,y_1) &= (100(-9),2(0.8))=(-900,1.6), \\
x_2 &= x_1-\eta(100x_1)= -9-0.1(-900)=81, \\
y_2 &= y_1-\eta(2y_1)=0.8-0.1(1.6)=0.64. &&
\end{flalign*}
می‌بینیم $x$ بین مقادیر بزرگ با علامت متناوب می‌پرد (نوسان شدید و واگرایی)، که همان مشکل بخش (ب) است.

\end{itemize}

\item \lr{AdaGrad}
\begin{itemize}
\item[(الف)]
در این تابع، گرادیان مؤلفه $x$ معمولاً بسیار بزرگ‌تر از $y$ است (به‌خاطر ضریب $100$ در مشتق).
\lr{AdaGrad} با جدا کردن نرخ یادگیری برای هر پارامتر، گام مؤثر $x$ را کوچک و گام مؤثر $y$ را بزرگ‌تر نگه می‌دارد؛ بنابراین با ناهمسانگردی (تفاوت انحنا/مقیاس) بهتر کنار می‌آید.

\item[(ب)]
در گام اول با $(x_0,y_0)=(1,1)$ داریم:
\begin{flalign*}
g_x^{(1)}=100,\qquad g_y^{(1)}=2. &&
\end{flalign*}
پس (با $\epsilon=10^{-8}$):
\begin{flalign*}
\eta_x^{(1)} &= \frac{0.1}{\sqrt{(100)^2}+10^{-8}}
\approx \frac{0.1}{100}=0.001,\\
\eta_y^{(1)} &= \frac{0.1}{\sqrt{(2)^2}+10^{-8}}
\approx \frac{0.1}{2}=0.05. &&
\end{flalign*}

\item[(ج)]
در \lr{AdaGrad} مخرج $\sqrt{\sum_{j=1}^k (g_i^{(j)})^2}$ با $k$ رشد می‌کند، بنابراین
$\eta_i^{(k)}$
به‌مرور کاهش می‌یابد و گام‌ها کوچک‌تر می‌شوند؛ اثر آن این است که الگوریتم می‌تواند در بلندمدت خیلی کند شود (به‌خصوص اگر گرادیان‌ها در ابتدا بزرگ باشند).

\end{itemize}

\item \lr{RMSProp}
\begin{itemize}
\item[(الف)]
\lr{AdaGrad} مجموع کل مربع گرادیان‌ها را جمع می‌کند و نرخ یادگیری را یکنواخت به سمت صفر می‌برد.
\lr{RMSProp} به‌جای جمع کل، از میانگین متحرک نمایی استفاده می‌کند:
\begin{flalign*}
s_{k+1}=\gamma s_k+(1-\gamma)g_k^2, &&
\end{flalign*}
بنابراین «حافظه محدود» دارد و نرخ یادگیری خیلی سریع به صفر میل نمی‌کند.
یک مزیت واضح: جلوگیری از کوچک شدنِ دائمی نرخ یادگیری و سازگاری بهتر با گرادیان‌های اخیر.

\item[(ب)]
در \lr{RMSProp} پارامتر $\gamma$ میزان حافظه/هموارسازی در تخمینِ مرتبه دوم (میانگین متحرک $g^2$) را کنترل می‌کند.
در \lr{Momentum} پارامتر $\beta$ میزان حافظه/هموارسازی در سرعت (میانگین متحرک گرادیان‌ها) را کنترل می‌کند؛ یعنی $\beta$ روی جهت حرکت (مرتبه اول) اثر می‌گذارد، ولی $\gamma$ روی مقیاس تطبیقی گام (مرتبه دوم) اثر دارد.
\end{itemize}

\item \lr{Adam}
\begin{itemize}
\item[(الف)]
در \lr{Adam}:
\begin{flalign*}
m_k &: \text{میانگین متحرک گرادیان }(\text{مرتبه اول، شبیه \lr{Momentum}}),\\
v_k &: \text{میانگین متحرک مربع گرادیان }(\text{مرتبه دوم، شبیه \lr{AdaGrad}/\lr{RMSProp}}).
&&
\end{flalign*}

\item[(ب)]
چون معمولاً $m_0=v_0=0$ مقداردهی اولیه می‌شوند، در چند گام اول $m_k$ و $v_k$ به سمت صفر بایاس دارند (کم‌برآورد می‌شوند).
تصحیح بایاس با تقسیم بر $(1-\beta_1^k)$ و $(1-\beta_2^k)$ این کم‌برآوردی را جبران می‌کند؛ بدون آن، در تکرارهای ابتدایی به‌روزرسانی‌ها بیش از حد کوچک می‌شوند و شروع یادگیری کند می‌شود.

\item[(ج)]
برای $f(x,y)=50x^2+y^2$ انتظار می‌رود \lr{Adam} بهتر از \lr{GD} ساده عمل کند چون:
\begin{itemize}
\item گام‌ها را به‌صورت جداگانه برای $x$ و $y$ مقیاس می‌کند (به‌خاطر $v_k$)، پس در راستای $x$ از جهش/نوسان می‌کاهد.
\item با $m_k$ (اثر \lr{Momentum}) نوسان‌های رفت‌وبرگشتی را هموار می‌کند و همگرایی را در دره‌های کشیده سریع‌تر می‌کند.
\end{itemize}
\end{itemize}
\end{itemize}